\documentclass{ltjsbook}
\usepackage{amsthm}
% 本文より前の設定は検査しない．
\newtheorem{Def}{定義}

\begin{document}

\chapter{群}\label{chap:group}

\section{定義}

\begin{Def}[群]\label{Def:group}
	集合\(G\)と写像\(\cdot \colon G \times G \to G\)の組\(\paren{G, \cdot}\)が%
	\term[ぐん]{群}であるとは，以下の条件をすべて満たすことをいう：
	\begin{enumerate}
		\item 任意の\(a, b, c \in G\)に対し\(\paren{a \cdot b} \cdot c = a \cdot \paren{b \cdot c}\)である．
		\item ある\(e \in G\)が存在し，任意の\(a \in G\)に対し\(a \cdot e = e \cdot a = a\)である．
	\end{enumerate}
\end{Def}

\begin{proof}
	\(b, b'\)がともに\(a\)の逆元であるとすると，
	\begin{equation*}
		b = b \cdot \paren{a \cdot b'} = \paren{b \cdot a} \cdot b' = b'
	\end{equation*}
	である．
	したがって逆元は一意である\footnote{単位元の一意性も同様に示せる．}．
\end{proof}

\begin{figure}
	\centering
	\includegraphics{group.pdf}
	\caption{群の例}
	\label{fig:group}
\end{figure}

\begin{verbatim}
x = 1
\end{verbatim}

\end{document}
